% Propietario conservado: /Users/ruben/Documents/New project/output/REVISION_ARGUMENTAL_APERTURAS_CIERRES_20260915/VII/delivery/MANUSCRIPT_VII_ES_EN_COMPLETE/source_es/manuscrito/29_retorno_y_memoria_traslacional.tex
% Artículo VII, recuperación aditiva, revisión 04, 10-09-2026.
% RESULTADO_RECUPERADO: S15, 11c_gravedad_holonomia_rev6.tex:16--126.
% El cociclo canónico c27=(1,18) se conserva como antecedente del cómputo
% de rutas; este fragmento prueba su composición y representación.
% La memoria traslacional finita y la amplitud cosmológica s0 tienen
% tipos distintos. El cociclo no evalúa por sí solo una densidad local.

\section{Phase return and representation of accumulated memory}
\label{vii:vii:sec:retorno-memoria-gravedad}

The chronology of enriched routes distinguishes position return,
orientation return and phase return. The observable publication
identifies certain arrival states, while the cumulative register
preserves the sequence traversed. The geometric representation must
transport both data jointly.

\subsection{Inversion cocycle and change of sheet}
\label{vii:vii:subsec:cociclo-memoria}

Let \(F_{\rm obs}\) be the chronological publication of TPK transport.
Its canonical blocks have the hierarchy
\[
\begin{array}{c|l}
27&\text{positional return with orientation reversal},\\
54&\text{positional and orientational return},\\
108&\text{complete return of the observable phase}.
\end{array}
\]
The register counts orientation inversions \(g(a)\) and
effective changes of sheet \(s(a)\). For a route,
\begin{equation}
 c(\gamma)=\sum_{a\subset\gamma}(g(a),s(a)),\qquad
 c(\gamma_2\gamma_1)=c(\gamma_2)+c(\gamma_1).
 \label{vii:vii:eq:cociclo-ruta}
\end{equation}
The canonical computation of the block of length \(27\) publishes
\(c_{27}=(1,18)\). Composition preserves those increments in
the four blocks that complete the observable period.
On forward-oriented routes the register is cumulative;
its extension to the route groupoid uses \(\mathbb Z^2\) and
\(c(\bar\gamma)=-c(\gamma)\).

\begin{proposition}[Memory of the complete observable turn]
\label{vii:vii:prop:memoria-vuelta-completa}
Writing \(\mathcal K_{27}=F_{\rm obs}^{27}\), the lift
to the register, for \(x\in\mathcal O_{108}\) and \(z\in\mathbb Z^2\), satisfies
\[
 \widetilde{\mathcal K}_{27}(x,z)
   =(\mathcal K_{27}x,z+(1,18)),\qquad
 \widetilde{\mathcal K}_{27}^{\,4}(x,z)
   =(x,z+(4,72)).
\]
After \(N\) complete turns, the increment is \(N(4,72)\).
\end{proposition}

\begin{proof}
The observable projection of \(\mathcal K_{27}\) has order four.
The preceding canonical computation fixes \(c_{27}=(1,18)\) as a constant
over that observable orbit. The additivity of
\eqref{vii:vii:eq:cociclo-ruta} sums the four increments and gives
\(c_{108}=4(1,18)=(4,72)\).
Repetition applies the same concatenation law and produces
\(Nc_{108}\). The return of the first coordinate preserves the
increment of the second.
\end{proof}

\subsection{Affine representation of the cocycle}
\label{vii:vii:subsec:representacion-afin-cociclo}

The discrete realization uses an element \(R_4\) of order four in
\(\operatorname{Spin}^+(1,3)\), a unit temporal vector \(u\)
fixed by its vector action and a unit of length \(\ell_0>0\).
The scale \(\ell_0\) retains its provenance in the dimensional section.
We define
\begin{equation}
 \rho_{\rm C}^{\rm disc}(\gamma)=
 \bigl(R_4^{c_g(\gamma)},\,\ell_0c_s(\gamma)u\bigr)
 \in\operatorname{Spin}^+(1,3)\ltimes\mathbb R^{1,3}.
 \label{vii:vii:eq:representacion-discreta-memoria}
\end{equation}
The power acts on the vector through the associated vector
representation of the spin group.

\begin{theorem}[Composition and translational amplitude of return]
\label{vii:vii:thm:traslacion-memoria}
The map \(\rho_{\rm C}^{\rm disc}\) preserves the identity,
concatenation and inversion. At complete returns,
\begin{equation}
 \rho_{\rm C}^{\rm disc}(\gamma_{108})=(I,72\ell_0u),\qquad
 \rho_{\rm C}^{\rm disc}(\gamma_{108}^{\,N})=(I,72N\ell_0u).
 \label{vii:vii:eq:traslacion-retorno-108}
\end{equation}
Inversion changes the sign of the translation.
\end{theorem}

\begin{proof}
The semidirect product is
\((R,a)(S,b)=(RS,a+Rb)\). Since \(R_4u=u\),
\[
\begin{split}
 \rho_{\rm C}^{\rm disc}(\gamma_2)
 \rho_{\rm C}^{\rm disc}(\gamma_1)
 &=
 \bigl(R_4^{c_g(\gamma_2)+c_g(\gamma_1)},
 \ell_0(c_s(\gamma_2)+c_s(\gamma_1))u\bigr)\\
 &=\rho_{\rm C}^{\rm disc}(\gamma_2\gamma_1).
\end{split}
\]
The identity cocycle is zero and that of the inverse is its opposite.
For the complete turn, substitute \(c_{108}=(4,72)\) and
\(R_4^4=I\). Concatenating \(N\) turns keeps the rotational part
equal to the identity and sums their translations.
\end{proof}

The amplitude \(72\ell_0\) is a length associated with a finite route.
A realization through a co-tetrad and a connection transports it to its
affine holonomy. The local contribution to the gravitational equations
is subsequently obtained through the curvature of that connection, the
variation of the matter action and the normalization of its density.
The representation thus preserves the exact cumulative content that
must enter that realization.
